\documentclass{article}
\usepackage{geometry}
\geometry{a4paper,scale=0.8}
\usepackage[utf8]{inputenc}

\usepackage{amsmath, amssymb}
\usepackage{amsthm}
\usepackage{mathrsfs}
\usepackage{enumitem}
\usepackage{blkarray}
\usepackage{tikz-cd}
\usepackage{hyperref}

\newtheorem{remark}{Remark}
\newtheoremstyle{breakstyle}
    {5pt}
    {10pt}
    {\normalfont}
    {0pt}
    {\bfseries}
    {\newline\indent}
    {0pt}
    {}
\theoremstyle{breakstyle}
\newtheorem{exercise}{Exercise}
% \let\ker\relax
% \DeclareMathOperator{\ker}{Ker}
% \DeclareMathOperator{\coker}{Coker}
% \DeclareMathOperator{\coim}{Coim}
% \DeclareMathOperator{\im}{Im}
% \let\dim\relax
% \DeclareMathOperator{\dim}{Dim}
% \let\hom\relax
% \DeclareMathOperator{\hom}{Hom}
% \let\max\relax
% \DeclareMathOperator{\max}{Max}
% \DeclareMathOperator{\rad}{Rad}
\DeclareMathOperator{\nil}{Nil}
% \DeclareMathOperator{\jab}{Jab}
% \DeclareMathOperator{\obj}{Obj}
% \DeclareMathOperator{\spec}{Spec}
\title{Chapter 3 - Rings and Moduls of Fractions}
\author{Flandre}
\date{22 Feb 2026}

\begin{document}
\maketitle

\begin{exercise}
    \leavevmode \vspace{-\baselineskip}
    \begin{enumerate}[label=({\roman*})]
        \item $(\Rightarrow)$:
            Let $M$ be generated by $\left\{ m_1,m_2,\ldots ,m_{n} \right\} $,
            for $\forall i\in [1,n]\cap \mathbb{Z}$, we have
            $$
            \forall r,t\in S,\exists s \in S,\text{s.t. }m_{i}s(r-t)=0,\text{i.e. }m_{i}s=0
            $$
            Denote the $s$ above as $s_{i}$.
            Since $M$ is finitely-generated, we can let
            $$
            s=s_1s_2\ldots s_n
            $$
            And have $sM=0$.
        \item $(\Leftarrow)$:
            When such $s$ exist,
            then for $\forall \frac{m_1}{s_1},\frac{m_2}{s_2}\in S^{-1}M$,
            $$
            s(s_1m_2-s_1m_2)=0\implies \frac{m_1}{s_1}=\frac{m_2}{s_1}
            $$
            So $S^{-1}M=0$.
    \end{enumerate}
    \qed
\end{exercise}

\begin{exercise}
    By \emph{Proposition 1.9},
    the problem is equivalent to prove that for any $\frac{a}{s}\in S^{-1}A$, we have
    $$
    \forall \frac{b}{t}\in S^{-1}A,1-\frac{a}{s}\frac{b}{t}\text{is a unit in }S^{-1}A
    $$
    Since $1+\mathfrak{a}$ is a multiplicative set, $\exists c\in \mathfrak{a},\text{s.t. }st=1+c$.
    Now
    \begin{align*}
        1-\frac{a}{s}\frac{b}{t}&= 1-\frac{ab}{st} \\
        &= 1-\frac{ab}{1+c} \\
        &= \frac{1+c-ab}{1+c}
    \end{align*}
    And we can easily see $\frac{1+c}{1+c-ab}=1+\frac{ab}{1+c-ab}\in S^{-1}A$.
    So the inverse of $1-\frac{a}{s}\frac{b}{t}$ is $1+\frac{ab}{1+c-ab}$.
    \qed
\end{exercise}

\begin{exercise}
    Define
    $$
    f:(ST)^{-1}A\longrightarrow U^{-1}(S^{-1}A)
    $$
    by mapping
    $$
    \frac{a}{st}\longmapsto \frac{\frac{a}{s}}{\frac{t}{1}}
    $$
    We'll prove $f$ is an isomorphism:
    \begin{enumerate}[label=({\roman*})]
        \item Surjection: 
            Trivial.
        \item Injection:
            $$
            \frac{\frac{a_1}{s_1}}{\frac{t_1}{1}}=\frac{\frac{a_2}{s_2}}{\frac{t_2}{1}}
            $$
            $\Leftrightarrow$ 
            $$
            \exists t\in T,\text{s.t. }\frac{t}{1}\left( \frac{a_1t_2}{s_1}-\frac{a_2t_1}{s_2} \right) =0
            $$
            $\Leftrightarrow$ 
            $$
            \exists t\in T,\text{s.t. }t(a_1t_2s_2-a_2t_1s_1)=0
            $$
            $\Leftrightarrow$
            $$
            \frac{a_1}{s_1t_1}=\frac{a_2}{s_2t_2}
            $$
            So $f$ is injective.
        \item Homomorphism:
            Trivial.
    \end{enumerate}
    \qed
\end{exercise}

\begin{exercise}
    Define
    $$
    f:S^{-1}B\longrightarrow T^{-1}B
    $$
    by
    $$
    \frac{b}{s}\longmapsto \frac{b}{f(s)}
    $$
    We'll prove it's an isomorphism:
    \begin{enumerate}[label=({\roman*})]
        \item Homomorphism:
            Trivial
        \item Surjection:
            Trivial
        \item Injection:
            $$
            \frac{b_1}{f(s_1)}=\frac{b_2}{f(s_2)}
            $$
            $\Leftrightarrow$ 
            $$
            \exists s\in S,\text{s.t. }f(s)\left( b_1f(s_2)-b_2f(s_1) \right) =0
            $$
            $\Leftrightarrow$ 
            $$
            \frac{b_1}{s_1}=\frac{b_2}{s_2}
            $$
            
    \end{enumerate}
    \qed
\end{exercise}

\begin{exercise}
    \leavevmode \vspace{-\baselineskip}
    \begin{enumerate}[label=({\roman*})]
        \item 
            By \emph{Proposition 3.11-(v)}, the operation $S^{-1}$ commutes with the formations of radicals.
            So $\forall \mathfrak{p}\subseteq A$,
            $$
            \left( \nil(A) \right) _{\mathfrak{p}}=\nil(A_{\mathfrak{p}})=0
            $$
            Therefore by \emph{Proposition 3.8}, $\nil(A_{\mathfrak{p}})=0$.
        \item 
            No.
            Counterexample given \href{https://math.stackexchange.com/questions/146951/if-the-localization-of-a-ring-r-at-every-prime-ideal-is-an-integral-domain-mu}{here}
    \end{enumerate}
\end{exercise}

\begin{exercise}
    \textbf{(UNFINISHED)}
    \begin{enumerate}[label=({\roman*})]
        \item 
            $\Sigma$ is nonempty, therefore has maximal elements by \emph{Zorn's Lemma}.
        \item 
            The prime property is trivial,
            we only need to show $A\backslash S$ is an ideal, and the minimal property can be obtained directly.
    \end{enumerate}
\end{exercise}

\begin{exercise}
    \leavevmode \vspace{-\baselineskip}
    \begin{enumerate}[label=({\roman*})]
        \item 
            \begin{enumerate}[label=({\Roman*})]
                \item $(\Leftarrow)$ :
                    When $A\backslash S$ is a union of prime ideals, $x\in A\backslash S$ or $y\in A\backslash S$ will immediately implies $xy\not\in S$.
                    So $xy\in S$ will imply $x,y\in S$.
                \item $(\Rightarrow)$ :
                    We'll prove that any elements in $A\backslash S$ was contained in a prime ideal which is disjoint to $S$:

                    $\forall a\in A\backslash S,(a)\cap S\subseteq \left\{ 0 \right\} $,
                    otherwise if $a^{n}=a\cdot a^{n-1}\in S$,
                    we'll have $a\in S$, contradictory.

                    Now we consider $\Sigma$ be the set of all ideals containing $(a)$ and being disjoint to $S$.
                    By \emph{Zorn's Lemma}, we can pick out a maximal element, namely $\mathfrak{p}\subseteq \Sigma$.
                    We'll prove $\mathfrak{p}$ is prime:

                    Suppose we have $xy\in \mathfrak{p},x,y\not\in \mathfrak{p}$.
                    Since $\mathfrak{p}$ is maximal and $\mathfrak{p}+(x),\mathfrak{p}+(y)\supsetneq \mathfrak{p}$,
                    we know both of $\mathfrak{p}+(x),\mathfrak{p}+(y)$ intersects with $S$.
                    Therefore $\exists s_1,s_1\in S$ such that $\exists p_1,p_2\in \mathfrak{p},a_1,a_2\in A$,
                    $$
                    s_1=p_1+a_1x,s_2=p_2+a_2x
                    $$
                    Then we have
                    $$
                    (s_1-p_1)s_2\in S
                    $$
                    whereas
                    \begin{align*}
                        (s_1-p_1)s_2&= a_1x(p_2+a_2y) \\
                        &= a_1p_2x+a_1a_2xy\in \mathfrak{p}
                    \end{align*}
                    Which means $(s_1-p_1)s_2\in S\cap \mathfrak{p}$, contradicting with the construction of $\mathfrak{p}$.
            \end{enumerate}
            \qed
        \item 
            By \emph{(i)}, $\overline{S}$ defined by the complement in $A$ of union of prime ideals disjoint from $S$ is indeed a saturated multiplicative closed subset of $A$.

            We only have to prove $\overline{S}$ is the unique minimal one,
            which is trivial since the family of sets, which consists of prime ideals disjoint from $S$, has a unique maximal element
            (final object, where morphisms taken under inclusion).
        \item 
            $$
            \mathfrak{p}\cap S\subsetneq \left\{ 0 \right\} 
            $$
            $\Leftrightarrow$ 
            $$
            \exists p\in \mathfrak{p},a\in \mathfrak{a},\text{s.t. }1+p=a
            $$
            $\Leftrightarrow$ 
            $$
            \mathfrak{a}+\mathfrak{p}=(1)
            $$
            $\Leftrightarrow$ 
            $$
            \mathfrak{a},\mathfrak{p}\text{ are coprime}
            $$
            Conversely, it's easy to prove that $\mathfrak{a},\mathfrak{p}$ being coprime will implies $\mathfrak{p}\cap S\subsetneq \left\{ 0 \right\} $.
            Therefore $\overline{S}$ is the complement of the union of all prime ideals which are not coprime with $\mathfrak{a}$.
    \end{enumerate}
\end{exercise}

\begin{exercise}
    \leavevmode \vspace{-\baselineskip}
    \begin{enumerate}[label=({\Roman*})]
        \item $(i)\Rightarrow(ii)$:
            $\forall t\in T,\phi(\frac{t}{1})=\frac{t}{1}$ is a unit in $T^{-1}A$.
            Since $\phi$ is bijective, the preimage of the inverse of $\frac{t}{1}\in T^{-1}A$ is also the inverse of $\frac{t}{1}\in S^{-1}A$.
        \item $(ii)\Rightarrow(iii)$:
            $\forall t\in T$, $\frac{t}{1}$ is a unit in $S^{-1}A$ and we denote $\frac{a}{s}=\left( \frac{t}{1} \right) ^{-1}$.
            Therefore $\exists s_0\in S,\text{s.t. }s_0(at-s)=0$,
            which indicate
            $$
            \left( s_0a \right) t=s_0s \in S
            $$
            Thus we can let $x=s_0a$.
        \item $(iii)\Rightarrow(iv)$:
            Suppose $\exists t\int T,\text{s.t. }t$ is contained in a prime ideal $\mathfrak{p}$ which is disjoint from $S$.
            By \emph{(iii)}, $\exists x\in A,\text{s.t. }xt\in S$.
            However, $t\in \mathfrak{p}\Rightarrow xt\in \mathfrak{p}\Rightarrow xt \in S\cap \mathfrak{p}$, contradictory.
        \item $(iv)\Rightarrow (v)$:
            Otherwise $T$ cannot be contianed in $\overline{S}$.
        \item $(v)\Rightarrow (i)$:
            
    \end{enumerate}
\end{exercise}

\end{document}
